% SOURCE: related_work.md §3.1–§3.5 (stubs), §4 (concessions to fold in),
% §6 (the only bibliography this section may draw from).
% Citation gate (broadened 2026-09-03, Decision Log #28): an identifier may
% be a Crossref DOI, an arXiv ID, or a PMLR / JMLR / ACL / DBLP record.
% Arya et al. 2019 (arXiv:1909.03012) and Mozannar & Sontag 2020 (PMLR v119)
% are cited here under that gate; both identifiers were resolved in-session.
% The §7 quarantine papers are still NOT cited. See the refs.bib header.
\section{Related Work}
\label{sec:related}

This work sits at the meeting point of two mature literatures, each of which
has settled a question we depend on and neither of which has asked the
question we ask. Each subsection below states what is already occupied.

\subsection{Stakeholder-tailored explanation}
\label{sec:rw-stakeholder}

That different audiences need different explanations of the same model
output is established, not open. Arya et al.\ made the point in a title ---
\emph{One Explanation Does Not Fit All} --- and built the toolkit and
taxonomy to go with it~\citep{arya2019explanation}. Imrie et al.\ show that
stakeholders in healthcare impose genuinely divergent interpretability
requirements, not differing degrees of one
requirement~\citep{imrie2023multiple}. Kim et al.\ build and compare
stakeholder-tailored explanation interfaces directly, including the
clinician-versus-patient split that two of our three tiers
reproduce~\citep{kim2024stakeholder}. Bello et al.\ go further and propose a
\emph{three-level} pipeline from technical explanation to natural language,
which is our structural twin~\citep{bello2025threelevel}. \textbf{We
therefore claim no novelty for the idea that explanations should be
audience-specific, nor for tiering as an architecture.} Both predate us.

Our contrast with Bello et al.\ is one of mechanism rather than structure,
and it is one of our stated contributions. Their mediator is a large language
model. Ours is a fixed set of templates over SHAP attributions, so every
string a health worker or a mother receives is determined by the case and can
be regenerated and audited. Nothing is generated at inference time and there
is nothing to hallucinate. We give up fluency and coverage for a rendering an
auditor can check line by line.

Suresh et al.\ present the sharpest challenge to how we have organised this
paper. They argue for moving \emph{beyond} role labels, since a stakeholder's
role determines neither the knowledge they bring nor the interpretability
task they face, and they offer a typology that decouples the
two~\citep{suresh2021beyond}. Our tiers are named by role: clinician, ASHA,
mother-to-be. We concede the point rather than wait for it to be made. Those
names are shorthand, and what actually varies across the tiers is what each
reader lacks, in reading modality, literacy and decision authority. Each tier
is designed against that deficit. In this deployment the two track each other
closely, since the ASHA is a specific, minimally-schooled cadre and the
mother may not read at all. Szymanski et al.\ report that role and expertise
co-vary in practice while remaining analytically
separable~\citep{szymanski2025disentangling}, which supports role as a
serviceable proxy here without rescuing it as a general axis. Where the two
come apart, a deployment would have to select on the deficit.

\subsection{Uncertainty communication, abstention, and selective prediction}
\label{sec:rw-uncertainty}

The second literature is where our claim lives, so we are precise about its
boundary. Withholding a prediction when the model is not confident is a
long-settled mechanism. Hendrickx et al.\ survey machine learning with a
reject option and name the case we implement: rejection under
\emph{ambiguity}, where two classes are close, as distinct from rejection
under novelty~\citep{hendrickx2024reject}. Mozannar and Sontag give the
learning-to-defer version, where the model is trained to route a case to a
human expert rather than answer it~\citep{mozannar2020consistent}. Our rule
declines the highest-severity signal when the top-two class probabilities are
separated by less than a fixed margin, which is a hand-set instance of the
same mechanism and a cruder one than either of these. \textbf{We claim none
of it.} Kompa et al.\ supply the clinical framing: in medical machine
learning, abstaining on high-uncertainty cases is a safety property rather
than a shortfall in coverage~\citep{kompa2021second}. That is the argument we
borrow to justify derating a signal instead of issuing it.

Bhatt et al.\ are the nearest ancestor of our headline
claim~\citep{bhatt2021uncertainty}. They treat uncertainty as itself a form
of transparency and consider explicitly how it should be \emph{communicated
to different stakeholders}. Our delta is narrow: they make the \emph{display}
of uncertainty audience-aware, whereas we make the \emph{suppression} of an
assertion audience-aware. The mother does not receive the uncertainty
rendered more gently, she receives a different assertion, while the clinician
receives the split probabilities unmodified. The abstention decision is taken
once, from the model's own output, then rendered per audience and graded so
that the reader with the least power to interrogate the system receives the
most conservative signal. We find no prior work that makes the abstention
\emph{rendering} a function of stakeholder vulnerability. That intersection,
and only that intersection, is what we claim.

\subsection{Explainable AI for community health workers in the Global South}
\label{sec:rw-chw}

Okolo et al.'s systematic review of explainable AI in the Global South
records how little of this work reaches the people it describes: across the
corpus they surveyed, deployment with actual users is close to
absent~\citep{okolo2022makingai}. That is the gap this paper sits in.

Three studies define our context and our principal exposure. Okolo et al.\
interviewed community health workers in rural India about AI-enabled mobile
health tools and about explanations
specifically~\citep{okolo2021cannot,okolo2024easy}. Solano-Kamaiko et al.\
built and studied an explorable explanation interface with the same
cadre~\citep{solanokamaiko2024explorable}. Srinidhi et al.\ document ASHA
Kirana, a deployment in which ASHAs collect data that feeds an algorithm and
returns a report for a doctor~\citep{srinidhi2021asha}, which is our data
flow exactly, minus the tiering and the uncertainty gate.

Two concessions follow, and both weaken us. First, this literature does not
report that simplification suffices. Okolo et al.\ and Solano-Kamaiko et al.\
both find community health workers struggling to interpret AI explanations
even when those explanations have been
simplified~\citep{okolo2024easy,solanokamaiko2024explorable}. The
comprehensibility of our ASHA card is therefore an open empirical question,
and that finding is the clearest statement of why a user study is the
necessary next step rather than an optional one. Second, every paper named in
this subsection has users and this paper has none. We propose a rendering
rule and evaluate the artifact that implements it, which is evidence of a
weaker kind than theirs, and not a claim to know better than they do what a
community health worker needs.

\subsection{Communicating risk without numeracy}
\label{sec:rw-riskcomm}

Our third tier renders to a reader who may not read, and the
risk-communication literature is both its warrant and its sharpest critic.
Galesic et al.\ show that icon arrays let low-numeracy readers extract risk
information they otherwise miss~\citep{galesic2009icon}, and
Garc\'{i}a-Retamero and Cokely's systematic review finds visual aids improve
risk literacy most among exactly the vulnerable readers our tier
targets~\citep{garciaretamero2017visual}. Richter et al.\ add the negative
result that matters most for a voice channel: purely verbal communication of
benefits and harms performs poorly for patients with limited health
literacy~\citep{richter2023communication}, which is why our spoken message is
paired with a visual rather than shipped alone. Peters et al.\ collect
current best practice for numeric risk~\citep{peters2025numeric}, and Bradley
et al.\ test risk presentation in the maternity
setting~\citep{bradley2025birth}.

That literature recommends conveying \emph{magnitude} transparently, through
an icon array or a visible denominator. Our mother tier strips magnitude
entirely and renders an ordinal band plus a single action. Read flatly, we
are doing what the field advises against.

The resolution turns on a precondition those recommendations assume: an icon
array communicates a magnitude the sender is entitled to assert. Our headline
case is precisely the one where the model is not so entitled, at a top-two
probability split of $0.487$ against $0.499$. Rendering that as roughly fifty
chances in a hundred would convey a precision the model does not possess, to
the reader least equipped to discount it. The departure is therefore
justified \emph{by} this literature rather than in spite of it: where the
magnitude is trustworthy, show it; where the uncertainty gate has fired, an
ordinal band plus an action is the honest rendering. That is
Section~\ref{sec:rw-uncertainty}'s argument arriving from the
health-communication side, and it remains an argument rather than a result,
since whether the band is understood as intended is untested here.

\subsection{Machine learning on the UCI Maternal Health Risk dataset}
\label{sec:rw-dataset}

Supervised classification with SHAP attributions on this dataset is
saturated. Mamun et al.~\citep{mamun2025identification}, Maisoon et
al.~\citep{maisoon2026stratification}, Widyawati et
al.~\citep{widyawati2025ensemble}, Rahman et al.~\citep{rahman2023xai},
Maheswari et al.~\citep{maheswari2024smote} and Hossen et
al.~\citep{hossen2024driven} all report the same broad construction: tree
ensembles, resampling for class imbalance, cross-validation, and SHAP for
post-hoc attribution. \textbf{There is no machine learning contribution
available here and we do not claim one.} Our model is a vehicle for the
rendering question, and that framing is forced by the field rather than
chosen out of modesty.

Maisoon et al.\ are our closest methodological neighbour and deserve explicit
separation~\citep{maisoon2026stratification}. They use the same dataset with
the same two-row outlier removal, tree ensembles with SMOTE, stratified
five-fold cross-validation, SHAP, and the same low-resource framing.
Everything upstream of the explanation is substantially their pipeline. What
is ours lies downstream: deduplicating before splitting and reporting what
that costs (Section~\ref{sec:data}), gating the rendering on the model's own
uncertainty, and rendering the gated output three ways. Their attributions
and Widyawati et al.'s independently rank blood sugar and systolic blood
pressure at the top~\citep{maisoon2026stratification,widyawati2025ensemble},
as ours do, which corroborates the attributions rather than adding a finding
of ours.
